\section{Mục tiêu trọng tâm trong tuần}
\begin{itemize}
    \item Nghiên cứu hệ thống tài liệu đặc tả kiến trúc phân hệ Trợ lý AI và giải pháp khai thác tri thức tài liệu.
    \item Nghiên cứu và thử nghiệm giải pháp RAG tích hợp cơ sở tri thức từ bộ tài liệu đặc tả kiến trúc phân hệ Trợ lý AI.
\end{itemize}

\section{Các công việc hoàn thành}

\subsection{Nghiên cứu Tài liệu Thiết kế Kiến trúc Phân hệ Trợ lý AI và Khai thác Tri thức}

\subsubsection{Mô tả chức năng}
Đọc và nghiên cứu toàn diện bộ tài liệu đặc tả kiến trúc, tiêu chuẩn tích hợp và kế hoạch triển khai của phân hệ Trợ lý AI:
\begin{itemize}
    \item \textbf{Tìm hiểu quy trình đặc tả tính năng:} Nghiên cứu tài liệu khảo sát kỹ thuật, tôn chỉ thiết kế và các trụ cột vận hành gồm nhận thức ngữ cảnh màn hình, cơ chế kiểm soát an toàn khi thực thi tác vụ và vòng lặp suy luận tự động có sự giám sát của người dùng.
    \item \textbf{Khảo sát kiến trúc cắm ghép và giao diện hội thoại:} Nắm bắt mô hình mở rộng qua các khe cắm giao diện chuẩn, sổ đăng ký các thành phần thẻ hội thoại và phương án tiếp nhận các mô-đun mã nguồn với giới hạn dung lượng tải nạp tối ưu.
    \item \textbf{Nghiên cứu giải pháp khai thác tri thức tài liệu:} Tiếp thu thiết kế tác nhân trích xuất thông tin tự thích ứng, mô hình kết hợp tìm kiếm ngữ nghĩa với đồ thị tri thức, kỹ thuật phân đoạn tài liệu bám cấu trúc và cơ chế đối chiếu trích dẫn trực tiếp trên trang văn bản gốc.
\end{itemize}

\subsubsection{Chi tiết kỹ thuật}
\begin{itemize}
    \item \textbf{Nắm vững các quyết định kiến trúc cốt lõi:} Tìm hiểu quy chuẩn làm phẳng dữ liệu ngữ cảnh giao diện, phương thức truyền dữ liệu thời gian thực dạng luồng, cơ chế điều hướng an toàn và nguyên tắc phân tầng cô lập lỗi độc lập.
    \item \textbf{Phân tích kịch bản kiểm thử và các lưu ý kỹ thuật:} Tiếp thu kinh nghiệm từ ma trận kiểm thử và danh mục các tình huống dễ phát sinh lỗi nhằm phòng ngừa rủi ro mất kết nối kênh truyền, xung đột trạng thái hoặc gián đoạn luồng xử lý.
    \item \textbf{Định hướng tích hợp thực tế:} Chuẩn bị cơ sở kỹ thuật phục vụ việc hoàn thiện khả năng duy trì ngữ cảnh trao đổi và hiển thị dẫn chứng tài liệu chính xác trên hệ thống.
\end{itemize}

\subsection{Ứng dụng Thực nghiệm Bộ Tri thức Kiến trúc vào Hệ thống RAG}

\subsubsection{Mô tả chức năng}
Tích hợp và đánh giá thực nghiệm việc ứng dụng bộ tài liệu đặc tả kiến trúc phân hệ Trợ lý AI vào hệ thống RAG thử nghiệm nội bộ, tập trung so sánh hiệu quả giữa mô hình RAG tiêu chuẩn và RAG nâng cao theo ngữ cảnh kiến trúc:
\begin{itemize}
    \item \textbf{Tích hợp cơ sở tri thức kiến trúc:} Chuyển hóa toàn bộ tài liệu đặc tả kiến trúc, sơ đồ luồng dữ liệu, quy chuẩn cắm ghép mô-đun và ma trận kiểm thử vào hệ thống RAG, tạo lập nguồn tra cứu kỹ thuật tập trung hỗ trợ hỏi đáp bằng ngôn ngữ tự nhiên.
    \item \textbf{Đảm bảo tính chính xác và đối chứng dẫn chứng:} Thiết lập cơ chế kiểm chứng câu trả lời luôn gắn liền với số trang và đoạn trích dẫn cụ thể trong tài liệu gốc; áp dụng nguyên tắc từ chối suy diễn khi thiếu dữ liệu nhằm loại bỏ hoàn toàn hiện tượng tạo thông tin sai lệch ngoài tài liệu.
    \item \textbf{Hỗ trợ cơ chế suy luận linh hoạt hai cấp độ:} Thử nghiệm chế độ phản hồi nhanh đối với các truy vấn tra cứu định nghĩa, tham số, mã lỗi và chế độ suy luận sâu cho các câu hỏi phức tạp cần đối chiếu chéo các ràng buộc kỹ thuật giữa nhiều phân hệ trước khi tổng hợp câu trả lời.
\end{itemize}

\subsubsection{Chi tiết kỹ thuật}
\begin{itemize}
    \item \textbf{Cơ chế nạp đa định dạng và phân đoạn phân cấp Cha - Con:} Tận dụng bộ phân đoạn chuyên dụng có khả năng bảo toàn cấu trúc bảng biểu thông số kỹ thuật và duy trì liên kết phân cấp giữa đề mục lớn và nội dung chi tiết, giúp giữ trọn vẹn ngữ cảnh kỹ thuật khi truy xuất.
    \item \textbf{Quản lý dữ liệu véc-tơ và thông tin ngữ cảnh:} Lưu trữ biểu diễn ngữ nghĩa của tài liệu vào cơ sở dữ liệu véc-tơ ChromaDB (hỗ trợ mô hình chạy nội bộ trên máy trạm hoặc dịch vụ ngoài), kết hợp cơ sở dữ liệu SQLite quản lý siêu dữ liệu về số hiệu phiên bản, đề mục và vị trí trang tài liệu.
    \item \textbf{Truy xuất lai kết hợp tái xếp hạng:} Phối hợp giữa tìm kiếm từ khóa chính xác và tìm kiếm ngữ nghĩa theo véc-tơ, áp dụng thuật toán tái xếp hạng để ưu tiên đúng các đoạn tài liệu có độ khớp cao nhất về thuật ngữ chuyên môn trước khi đưa vào ngữ cảnh cho mô hình ngôn ngữ.
    \item \textbf{Giao diện tương tác và kiểm chứng trực quan:} Triển khai giao diện thử nghiệm hỗ trợ hiển thị phản hồi thời gian thực, trực quan hóa các bước suy luận trung gian và gắn nhãn dẫn chứng nguồn giúp người dùng nhấp xem trực tiếp trang tài liệu đối chứng.
\end{itemize}

\section{Các công việc đang thực hiện}
\begin{itemize}
    \item Rà soát lỗi, tìm hiểu và áp dụng thử các tài liệu liên quan kiến trúc hệ thống, chuẩn bị kế hoạch tiếp theo.
\end{itemize}
